% 论文入口。沿用原 cumcmthesis.cls / cumcm2026.sty；原 example.tex 未覆盖。
% 建议用户在本目录用 XeLaTeX 编译两次。本次仅提供源稿，不编译。
\documentclass[withoutpreface,bwprint]{cumcmthesis}
\usepackage{cumcm2026}
\usepackage{url}
\usepackage{placeins}
\hypersetup{hidelinks}
\setlength{\emergencystretch}{2em}
\renewcommand{\arraystretch}{1.13}
\title{预测误差驱动的微网购电与储能协同调度}
\tihao{C}
\yearinput{2026}
\monthinput{9}
\dayinput{10}

\begin{document}
\maketitle
\begin{abstract}
针对光伏微网日前购电、储能控制和预测更新相互影响的问题，建立由确定性调度、残差场景优化、定时滚动反馈和波动电价联合场景逐步扩展的统一模型。以10分钟为基本时段，采用电网侧电量描述充放电，分别计入两侧90\%效率，并连续传递日际储电量。通过历史误差构造参考备用和紧急费用风险项，在电费与应急购电需求之间形成可解释的权衡。

问题一建立购电与储能联合线性规划，并利用免费弃余条件证明存在不同时充放电的最优解。计算得到全天计划购电59482.70 kWh、购电费35126.95元，较无储能基准节省26.90\%，首末储电量均为6000 kWh。

问题二以1月进行预测模型选择，以2月至12月334天进行顺序回测。负荷采用最近四个同星期日均值，光伏采用最近七日均值。利用此前28日整条残差，经代表场景缩减构建共享储能参考轨迹的两阶段随机规划，并加入紧急费CVaR与参考备用。实际总费用为1405.86万元，较点预测LP降低11.69\%，紧急购电量为79905.90 kWh。

问题三在0、6、12、18时更新剩余日计划，按原始承诺与最终生效计划计算非对称调整费。总费用为1404.27万元，紧急购电量为12685.74 kWh。与采用同一零时预报而全天不更新的策略相比，费用降低1.19\%，紧急电量降低83.03\%。更新收益包含预报、实测状态及负荷偏差反馈的共同作用。

问题四仅以历史及当前可见数据预测电价，通过同日联合残差保留价格与净负荷关联，得到静态与滚动策略费用分别为1487.45万元和1487.20万元。滚动策略的费用改善有限，但紧急电量从75130.77降至9795.55 kWh。敏感性分析表明，增加备用或风险权重通常以额外采购换取较少紧急需求，不能保证实际总费用尾部同时改善。本文据此区分经济最优偏好与应急需求控制偏好，并给出完整策略、计算代码及指定日期结果。

\keywords{微网调度\quad 两阶段随机规划\quad 残差场景\quad 滚动优化\quad 储能备用}
\end{abstract}

\section{问题重述与分析}

光伏发电具有日内波动，居民负载与光伏出力的高峰也未必重合。并网微网可利用储能转移电量，并根据外网电价确定购电时段。其调控任务是在满足负载和储能物理约束的条件下，协调提前购电、储能运行与事后补购，降低实际支付费用。储能既可用于峰谷价差套利，也可用于吸收预测误差，两种用途竞争同一有限容量。

问题一给定单日电价、负载和光伏预测，要求每天零时形成购电计划，并使日首、日末储电量相同。该问题构成后续模型的物理基础：以时段购电量和储能充放电量为控制变量，以储电量为跨时段状态。日前储能调度及其线性建模可参考文献\cite{kong2016,silva2020}，本题的设备参数和约束数值均取自赛题。

问题二中负载和光伏实际值在决策时尚未知晓。计划买少会触发五倍电价的紧急购电，买多又不能获得退款，因此不能仅对预测均值重复求解问题一。本文用历史预测残差构造场景，将计划决策与紧急补救分层，并对高额紧急费用加入风险度量。两阶段随机调度与风险约束研究\cite{abunima2022,marino2018,grover2018}提供了这一结构的依据，但本题不引入文献中的柴油机、风电机组、可中断负载和售电收入。

问题三在每日0:00、6:00、12:00、18:00获得新光伏预报，可据此修订未执行的购电量。预测更新与调整费用必须同时考虑：预报更准确，并不自动意味着支付调整费用后仍然更省钱。本文把问题二的优化器置于定时滚动框架内，用实际储电量连接前后两次决策\cite{palma2013,silvente2015,li2016,jani2022}。问题四进一步将电价视为未知的未来交易价格，在保留物理约束与结算方式的基础上加入价格预测及联合残差场景\cite{li2016,li2020}。

四问采用共同的“预测、计划、执行、反馈”结构，递进增加不确定性和信息更新。模型的针对性设计是利用近期预测误差同时调节紧急费用风险权重与储能参考备用，并将优化参考轨迹和实际可执行控制分开。该设计是否节省费用，应以附件数据上的实际结算与对照计算判断，不能仅凭模型形式认定其优越性。

\section{模型假设与符号说明}

\subsection{基本假设}

\begin{enumerate}
  \item 将输入功率样本解释为对应10分钟时段的平均功率，以$\Delta t=1/6$小时换算电量。忽略微网内部线路损耗和电池自放电，电池充电、放电效率分别为$\eta_c=\eta_d=0.9$，故完整往返效率为0.81。
  \item 不向外网售电；超过负载与可充电能力的富余电量允许免费弃置。计划购电按承诺电量付费，未使用电量不退款；不额外引入题目未给出的寿命折损、碳排放或设备投资费用。
  \item 储能容量及功率限制按附录执行。问题一采用$e_0=e_{144}=6000$ kWh；问题二至四只在2025年1月1日零时初始化为6000 kWh，以后传递前一日实际末态，不每日重置。
  \item 外网可在当前时段补足剩余缺口，题目未给定紧急购电容量上限，故不另设此上限。由此实际负载能够得到满足，需优化的是补购费用与紧急购电发生情况，而不是宣称提高某一失供电概率指标。
  \item 在10分钟执行尺度上，可观测当前时段负载、光伏与储电量并即时平衡功率。未来实际值不可提前使用；历史价格、负载以及当前预报发布前已经结束时段的测量值均可用于决策。
  \item 附件4记录实际交易电价，不将当天全部价格提前告知日前优化器。预测和场景用于决策，实际价格用于结算，二者严格分离。
\end{enumerate}

\subsection{主要符号}

下文所有购电、充放电与储电变量均以kWh为单位，只有原始输入采用kW。为避免混淆，以$p$表示电价，以$c$表示充电量。

\begin{table}[htbp]
\centering
\caption{主要符号及含义}
\label{tab:symbols}
\small
\begin{tabular}{p{0.18\textwidth}p{0.62\textwidth}p{0.11\textwidth}}
\toprule
符号 & 含义 & 单位\\
\midrule
$D,t,s$ & 日期、日内时段、场景索引，$t=0,\ldots,143$ & ---\\
$L_{D,t},R_{D,t}$ & 实际负载功率与光伏发电功率 & kW\\
$n_{D,t}$ & 净负荷电量，$(L_{D,t}-R_{D,t})\Delta t$ & kWh\\
$p_{D,t}$ & 当前时段实际交易电价 & 元/kWh\\
$g_{D,t},a_{D,t}$ & 零时计划购电量、最终执行的调整购电量 & kWh\\
$c_{D,t},d_{D,t}$ & 电网侧充电输入量、放电输出量 & kWh\\
$e_{D,t}$ & 时段$t$开始时的储电量 & kWh\\
$h_{D,t},w_{D,t}$ & 紧急购电量、富余弃置电量 & kWh\\
$\pi_s,Q_s$ & 场景概率、场景整段紧急购电费用 & ---，元\\
$\rho_t,\lambda$ & 参考储能备用量、风险权重 & kWh，---\\
$E_{\min},E_{\max}$ & 储电量下限1200、上限10800 & kWh\\
$U$ & 单时段充放电量上限，$5000/6$ & kWh\\
\bottomrule
\end{tabular}
\end{table}

\section{数据处理与因果预测}

\subsection{时间尺度与结果表头约定}

附件2和附件4包含2025年365天、每天144个时刻的数据，附件3包含每天四次发布的24点整点光伏预报。对日期顺序、样本数、有限值和非负性逐项检查，原始负载、光伏、电价及预报均未发现缺失或负值；保留观测极值，不用全年统计量事后删除异常值。原始文件只读保存，并记录文件哈希以支持追溯。

源数据中的00:10作为右端点，对应物理区间00:00--00:10；最后一个标记$0{:}00{+}1$对应当日23:50--24:00。于是
\begin{equation}
 I_t=[t/6,(t+1)/6),\qquad
 n_{D,t}=(L_{D,t}-R_{D,t})/6.
 \label{eq:time_energy}
\end{equation}
最后一个区间的右端点24:00归入前一日结算。全文指定时段表均按上述自然日物理区间取数，四小时充放电量由24个时段求和。

附件5原模板的首列区间显示为0:10--0:20，末列显示为$0{:}00$--$0{:}10{+}1$，与上述物理区间存在10分钟显示偏移。按用户指定保留原表头及其格式，结果按源数据144列的顺序逐列一一填入：第一个数据列保存$I_0$结果，最后一列保存$I_{143}$结果。该规则是列序映射，不按模板区间字面再平移一次数据；工作簿显示标签不改变优化器的自然日边界。

\subsection{训练期划分与预测模型选择}

2025年1月用于建立预测、积累历史误差并推进储电状态；2月1日至12月31日的334天作为结果填报期。候选预测器包括前一日、3/7/14/28日移动平均、前一周、最近2/4个同星期日均值、同星期均值与近7日均值的加权组合，以及带近期水平修正的两种预测，共12种。对每个候选器，预测日期$D$时只能读取$D$以前的实际数据。

模型类型依据1月15日至31日逐日前推预测的均方根误差选择，2月以后固定类型，仅更新其可用历史窗口。负载和电价选中最近至多4个同星期日均值，历史光伏预测选中近7日均值。以$X$代表负载或电价，其预测形式为
\begin{equation}
 \widehat X_{D,t}=\frac{1}{|\mathcal W_D|}\sum_{j\in\mathcal W_D}X_{j,t},
 \quad \mathcal W_D=\{D-7,D-14,D-21,D-28\}\cap\{1,\ldots,D-1\},
 \label{eq:weekday_forecast}
\end{equation}
光伏则平均此前至多7天同一时段出力。历史不足时使用已有样本；首日用附件1作为固定先验。1月前15天使用近7日均值，此后仅按已经发生的1月预测误差在线选型，不把1月底才确定的模型选择结果反用于月初决策。

均方根误差用于选型，平均绝对误差用于补充描述，两者分别为
\begin{equation}
 \mathrm{RMSE}=\sqrt{\frac1N\sum_{i=1}^{N}(X_i-\widehat X_i)^2},
 \qquad \mathrm{MAE}=\frac1N\sum_{i=1}^{N}|X_i-\widehat X_i|.
\end{equation}
全年误差统计只评价所选方法，不参与重新选型。储能在1月也逐时执行并向后传递，因此2月1日初态是训练期运行的实际结果。

\subsection{发布预报插值与日内修正}

问题二不使用附件3，光伏完全依据历史预测；问题三及问题四对应的滚动策略使用当前已发布预报。将发布时刻记为$r\in\{0,6,12,18\}$小时，24个预报点解释为$r+1$至$r+24$小时。在发布时刻已观测光伏值与这些预报点之间作线性插值，得到10分钟右端点的预测；零时锚点取前一日24:00观测，首日零时取0。仅保留发布以后且不超过本日24:00的部分，跨日预报在当前实现中不进入优化。

日内负载预测以零时预测为基线，用此前2小时已结束的12个时段计算平均偏差$b_r$。设$k=1,2,\ldots$为发布后的预测步数，采用
\begin{equation}
 \widehat X_{r,k}=\max\!\left\{0,\widehat X^{(0)}_{r,k}
 +\operatorname{clip}\!\left(b_r\exp(-k/36),-0.3\widehat X^{(0)}_{r,k},0.3\widehat X^{(0)}_{r,k}\right)\right\}.
 \label{eq:issue_correction}
\end{equation}
其中$\operatorname{clip}$表示限制到给定上下界。6小时衰减尺度和30\%修正幅度预先固定；波动电价按同一规则修正，并设置正的数值下界。用于计算$b_r$的都是发布前已观测值。光伏直接采用本次发布预报，不使用事后实际曲线替代。

\section{统一购电与储能调度模型}

\subsection{问题一：确定性日前联合优化}

先省略日期下标。在已知单日输入条件下，令$h_t=0$、$a_t=g_t$，以购电费最小为目标：
\begin{align}
 \min\quad &\sum_{t=0}^{143}p_tg_t,\label{eq:q1_objective}\\
 \text{s.t.}\quad &g_t+d_t-c_t\ge n_t,\label{eq:balance}\\
 &e_{t+1}=e_t+\eta_cc_t-\frac{d_t}{\eta_d},\label{eq:soc}\\
 &E_{\min}\le e_{t+1}\le E_{\max},\quad 0\le c_t,d_t\le U,\label{eq:limits}\\
 &g_t\ge0,\qquad e_0=e_{144}=6000.\label{eq:terminal_q1}
\end{align}
式\eqref{eq:balance}中未消纳的非负差值就是弃置电量$w_t$，也可显式写成供需等式。状态方程中的$c_t$为电池外部输入电量，$d_t$为电池外部输出电量，因此分别乘充电效率、除放电效率，而不是把两者都乘0.9。

该模型为线性规划。由于允许免费弃置，同步充放电无需增加二进制变量。若某时段$c_t,d_t>0$，令
\begin{equation}
 \varepsilon=\min\{c_t,d_t/(\eta_c\eta_d)\},\qquad
 c'_t=c_t-\varepsilon,\quad d'_t=d_t-\eta_c\eta_d\varepsilon.
\end{equation}
此变换保持储电量变化不变，使充电净需求减少$(1-\eta_c\eta_d)\varepsilon$，因而不破坏供电约束，也不增加购电费，并消去至少一种充放电动作。存在满足互斥的最优解。程序给充放电总量加入$10^{-7}$元/kWh的数值破同解项以优先选取此类解，实际费用统计不计入该项。

忽略容量和功率的机会成本时，若在低价$p_{\ell}$时段购入电量，并在高价$p_h$时段放电替代购电，盈利的必要价差条件为
\begin{equation}
 p_h>\frac{p_{\ell}}{\eta_c\eta_d}=\frac{p_{\ell}}{0.81}.
 \label{eq:arbitrage}
\end{equation}
实际调度还受充放电功率、剩余容量、光伏剩余量和日末回补约束影响，不能仅根据相邻电价大小作局部选择。

\subsection{历史整日残差与场景缩减}

在日期$D$、决策时刻$\tau$，历史集合$\mathcal H_D$取此前至多28天。为避免固定先验首日残差长期参与场景估计，历史充足后从1月2日起取样；仅首日无历史时使用单点预测。相同发布时间的历史净负荷误差为
\begin{equation}
 \epsilon^{n}_{j,t}
 =\left[(L_{j,t}-R_{j,t})-(\widehat L_{j,t}^{(\tau)}-\widehat R_{j,t}^{(\tau)})\right]\Delta t,
 \quad j<D,\quad t\ge\tau.
 \label{eq:residual}
\end{equation}
问题二用历史零时预测误差，滚动问题用相同发布时刻的预测误差。一个历史日构成一条完整剩余日误差轨迹，不逐点独立抽样，也不强制假定正态分布。

以轨迹间欧氏距离作确定性$k$-medoids缩减，最多保留12个历史代表日。先选总距离最小的轨迹，再逐次选距离已有代表最远的轨迹；交替分配与更新中心，最多8轮。代表日仍是实际存在的历史误差轨迹，场景概率为其所属簇样本数占历史样本数的比例：
\begin{equation}
 n_{t,s}=\widehat n_t+\epsilon^n_{j_s,t},\qquad
 \pi_s=\frac{|\mathcal C_s|}{|\mathcal H_D|},\qquad \sum_s\pi_s=1.
\end{equation}
历史不足12条时全部保留，赋予相等概率。该方法只对近28日经验分布作压缩，其场景覆盖能力受短窗口限制，不表示已经获得未来随机过程的真实分布。

\subsection{问题二：受限两阶段随机规划}

日前阶段同时选择场景共享的购电$g_t$、参考充放电$c_t,d_t$及参考储电量$e_{t+1}$；第二阶段仅允许各场景的紧急购电$h_{t,s}$不同。对所有场景要求
\begin{equation}
 g_t+d_t+h_{t,s}-c_t\ge n_{t,s},\qquad h_{t,s}\ge0.
 \label{eq:scenario_balance}
\end{equation}
因此本模型属于限制策略类的两阶段随机规划：储能参考轨迹作为第一阶段共同决策，不是每个场景拥有一条完全自适应的储能轨迹。这样减少了变量数量，也避免在规划时把整日未来信息提前赋予储能控制，但可能牺牲更一般多阶段策略的经济性。实际执行不选择某个事后最有利场景，而采用下述当前状态反馈规则。

令$Q_s=\sum_{t=\tau}^{143}5p_{t,s}h_{t,s}$为场景在整个剩余时域内的紧急费用，$\bar p_t=\sum_s\pi_sp_{t,s}$。固定电价问题中$p_{t,s}=p_t$。采用期望费用、尾部紧急风险及日末储电偏离的目标：
\begin{equation}
 \min\left\{\sum_t\bar p_tg_t+\sum_s\pi_sQ_s
 +\lambda\operatorname{CVaR}_{0.90}(Q)
 +\kappa|e_{144}-6000|\right\}.
 \label{eq:q2_objective}
\end{equation}
约束包括式\eqref{eq:soc}、\eqref{eq:limits}和\eqref{eq:scenario_balance}，优化初态为当前实际储电量。CVaR衡量完整轨迹紧急费的尾部，而非将各时段CVaR简单相加，也不是全部购电费用的CVaR。依照文献\cite{rockafellar2000}，引入阈值$z$和超额变量$v_s$，写为
\begin{equation}
 \operatorname{CVaR}_{\alpha}(Q)
 =\min_z\left(z+\frac{1}{1-\alpha}\sum_s\pi_sv_s\right),
 \quad v_s\ge Q_s-z,\quad v_s\ge0,
 \quad\alpha=0.90.
 \label{eq:cvar}
\end{equation}
因$Q_s\ge0$，阈值限制为非负不损失本题的最优值。日末绝对值同样通过两个非负变量线性化。取$\kappa=\operatorname{median}_{t}(\bar p_t)$，以接近本时域典型电价的系数约束末端价值损失。该项是优化用软目标，不计入实际电费；问题二至四不强制每日回到6000 kWh，实际次日继续继承末态。

\subsection{预测误差驱动的风险与参考备用}

用近期净负荷误差相对尺度决定风险权重。设$\sigma_t$为历史$\epsilon^n_{j,t}$的标准差，$N_\tau=144-\tau$，定义
\begin{equation}
 r=\frac{N_\tau^{-1}\sum_{t=\tau}^{143}\sigma_t}
 {\max\{N_\tau^{-1}\sum_{t=\tau}^{143}|\widehat n_t|,50\}},
 \qquad \lambda=0.10\min\{3,r/0.20\}.
 \label{eq:risk_weight}
\end{equation}
分母中的50 kWh防止净负荷接近零时比值失稳，权重上限为0.30。误差越大，目标函数越关注少数高额紧急补购场景；该经验映射的作用强度受给定参数限制。

对时段$t$结束后的未来至多1小时，取6个时段净负荷误差累计值的80\%经验分位数，换算为电池内部应保留的参考电量：
\begin{equation}
 \rho_t^{\rm raw}=\min\!\left\{2400,
 \frac{\max\!\left[0,q_{0.8}\!\left(\sum_{k=t+1}^{\min(t+6,143)}\epsilon^n_{j,k}\right)\right]}{\eta_d}\right\}.
 \label{eq:reserve}
\end{equation}
累计误差沿用同一历史轨迹，保留时段间共同偏差；临近日末时只使用本日尚余时段，空和为零。若当前时刻为$\tau$、实测储电量为$e_\tau$，再按最大充电速度截断为可达备用：
\begin{equation}
 \rho_t=\min\{\rho_t^{\rm raw},
 \max[0,e_\tau-E_{\min}+\eta_cU(t-\tau+1)]\},\quad
 e_{t+1}\ge E_{\min}+\rho_t.
 \label{eq:reachable_reserve}
\end{equation}
备用上限2400 kWh约为可用容量的四分之一，不把全部电量固定为风险缓冲。CVaR从费用尾部施加偏好，备用从储能参考路径施加约束，两者处理的对象不同。

需要区分参考备用和真实硬约束：式\eqref{eq:reachable_reserve}约束优化参考轨迹，实际运行在当前缺口较大时可以消耗这部分电量，硬性保持的只有1200 kWh物理下限。由于外网紧急补购始终可用，该备用设计不构成概率供电保证；也不能把80\%经验分位解释为具有严格校准的80\%可靠性约束。

\subsection{因果执行器与状态传递}

计划形成后，实际负载与光伏可能偏离参考曲线。执行器在每个10分钟时段只接收当期有效购电$a_t$、当前净负荷$n_t$和当前储电量$e_t$。记$[x]_+=\max(x,0)$。若$a_t\ge n_t$，采用
\begin{equation}
 c_t^{\rm act}=\min\{a_t-n_t,U,(E_{\max}-e_t)/\eta_c\},\quad
 d_t^{\rm act}=h_t=0,\quad w_t=a_t-n_t-c_t^{\rm act}.
 \label{eq:execute_charge}
\end{equation}
若$a_t<n_t$，则
\begin{equation}
 d_t^{\rm act}=\min\{n_t-a_t,U,\eta_d(e_t-E_{\min})\},\quad
 c_t^{\rm act}=w_t=0,\quad h_t=n_t-a_t-d_t^{\rm act}.
 \label{eq:execute_discharge}
\end{equation}
之后由式\eqref{eq:soc}更新实际储电量。该规则把净缺口投影到当前可行充放电区间，再用紧急购电补足剩余缺口；多余电量则充电或弃置。实际充放电由这一反馈规则产生，不能将优化器内部的$c_t,d_t$直接当作全年实际执行量。

由两种情况可逐步验证，充放电互斥、储电量范围、功率上限及含紧急购电的供需平衡均满足。对问题二$a_t=g_t$，购电计划全天不改，但储能可根据当期观测实施反馈。日际满足
\begin{equation}
 e_{D+1,0}=e_{D,144},
 \qquad e_{2025\text{-}01\text{-}01,0}=6000.
\end{equation}
因此1月训练期的控制行为会影响2月初状态，不能在填报开始时重新初始化以获得额外免费电量。

\subsection{问题三：结算约束下的定时滚动优化}

在零时用当前发布预报确定原始计划$g_t$；在6:00、12:00、18:00，即$\tau=36,72,108$，根据新预报、历史同发布时间误差以及当前实际储电量，重新优化$t=\tau,\ldots,143$。已执行时段保持不变，每次仅执行直到下一个规定更新时间的购电量。由此形成递减时域的定时MPC；尽管新预报覆盖未来24小时，本实现每次只优化至本日24:00。

采用赛题文字的主结算口径：零时计划全额付费；下调另支付下调量对应的50\%违约费，上调量另按1.5倍电价支付；紧急购电另按5倍支付。令
\begin{equation}
 \delta_t^+=[a_t-g_t]_+,\quad\delta_t^-=[g_t-a_t]_+,
\end{equation}
则每天实际总费用为
\begin{equation}
 C_D=\sum_t p_{D,t}g_{D,t}
 +\sum_t p_{D,t}(1.5\delta_{D,t}^++0.5\delta_{D,t}^-)
 +\sum_t5p_{D,t}h_{D,t}.
 \label{eq:settlement}
\end{equation}
该口径未增加题目未给定的取消计划退款。修订始终与零时原始计划比较，避免把多次更新之间的差值重复收费。在每次优化中使用$a_t-g_t=u_t-v_t$、$u_t,v_t\ge0$线性化调整量，并将$1.5\bar p_tu_t+0.5\bar p_tv_t$加入目标；已经承诺的计划费用相对于当前决策为常数，可从当前优化目标中略去。紧急费用、CVaR和日末软目标仍保留。

\paragraph{下调购电的支配性。}
在正电价、计划不退款、弃余免费的假设下，若某时段$a_t<g_t$，把$a_t$提高为$g_t$，保持原有充放电和紧急购电不变，将新增电量弃置，仍满足全部物理约束，却免去了$0.5p_t(g_t-a_t)$违约费。因此$a_t<g_t$受到支配，至少存在一个满足$a_t\ge g_t$的最优修订解。该性质来自本题结算和弃置假设，不是求解器禁止下调。它也说明，新的预报主要通过发现购电不足和调整储能参考路径发挥作用，光伏上调预报未必能收回原计划费用。

上述结论针对优化可行域和目标函数；实际控制器还根据当前储电量执行反馈，不能因此推断所有更新必然降低实现费用。是否需要额外时刻的预报，须在相同结算方式下比较获得信息后的可执行费用收益与相应成本；现有附件仅支持四个规定发布时刻，不据此编造更高更新频率的实测效果。

\subsection{问题四：价格与净负荷联合不确定性}

波动价格预测使用式\eqref{eq:weekday_forecast}，滚动时再按式\eqref{eq:issue_correction}用已观测价格修正。历史价格残差为$\epsilon^p_{j,t}=p_{j,t}-\widehat p_{j,t}^{(\tau)}$，与相同日期、相同发布时间的净负荷残差拼接成联合轨迹。净负荷残差本身由负载残差减光伏残差形成，因此不把价、荷、光伏随机误差分别打乱组合。

场景缩减前，对净负荷误差块和价格误差块分别除以各块的总体标准差，再计算欧氏距离，防止kWh数值尺度淹没价格误差。选中同一个代表日$j_s$后同时生成
\begin{equation}
 n_{t,s}=\widehat n_t+\epsilon^n_{j_s,t},\qquad
 p_{t,s}=\max\{0.01,\widehat p_t+\epsilon^p_{j_s,t}\}.
 \label{eq:joint_scenario}
\end{equation}
该构造保留样本中的净负荷与价格关联及日内轨迹关系，不声称建立了负载、光伏和电价的完整联合概率分布。0.01元/kWh仅是优化场景的保守支撑下界，用于避免加性残差生成非正电价后出现无限购电激励；它不是原始数据清洗阈值。附件实际最低价0.0076元/kWh予以保留，全部实际费用按未截断价格计算。

问题四对应问题二时，用联合场景在零时优化，全天购电计划不变；对应问题三时，在四个规定时刻更新并执行滚动策略。两者均以场景均价计入第一阶段购电或调整成本，以场景各自价格计算紧急费用和尾部风险，随后按式\eqref{eq:settlement}对实际价格结算。风险权重和参考备用仍由净负荷误差决定，价格风险通过联合场景和紧急费CVaR进入，不额外声称实现了价格驱动的自适应风险权重。

\subsection{求解流程与适用范围}

所有场景系数确定后，上述目标及约束均为线性形式，采用稀疏矩阵构建，由HiGHS求解线性规划。每个决策点依次执行：读取当前可用信息，更新预测，选择历史残差与代表场景，计算风险权重及可达备用，求解剩余日计划，按当前实际净负荷逐时执行，并把真实末态传至后续决策。零时保留原始计划副本供结算和结果填报，后续修订不覆盖这一基准。

该实现求解的是所设限制策略类内的优化问题，不是微网任意因果策略的全局最优控制。整日共享储能参考、至日末的有限优化时域、28日经验窗口以及固定反馈执行器，均可能带来经济性损失。终端软目标减少末端过度放电倾向，但不能替代跨日长期价值函数。本文分别报告优化中的风险目标与实际支付费用，模型预估紧急费不代替实际执行器产生的紧急费。

\section{模型求解与结果分析}

\subsection{结果口径与评价范围}

问题一使用附件1的一天数据。其余问题均从2025年1月1日以6000 kWh初态启动，1月用于历史预测选择、残差积累和实际储能状态推进；正式填报与策略比较覆盖2月1日至12月31日，共334天、48096个10分钟时段。各策略独立推进，因此2月1日的实际初态可以不同，不能把1月运行后的库存差异隐去。

论文表格采用物理自然日时段；附件5购电页严格保留原始时间标签和列顺序，源数据与结果逐列对应，具体映射见数据处理节及附录。指定时段购电量与全天常规购电量均不含紧急购电，紧急量另列。论文中的“全天购电费”统一为全部实际支付费用。工作簿计划购电页的费用列则对应计划费用，调整购电页的费用列对应全部实际费用，避免把常规计划费用和应急费用混为一项。

四组策略的总费用、紧急需求及状态边界见表\ref{tab:annual}。其中日费用CVaR95是334个实际日费用的经验尾部均值，用于描述高费用日期，不是置信区间，也不等同于优化目标中对场景紧急费使用的CVaR90。固定电价与波动电价属于不同结算环境，不能把两者总费用之差全部归因于算法。

\begin{table}[htbp]\centering\small\caption{2025年2月1日至12月31日费用与风险结果}\label{tab:annual}
\begin{tabularx}{\linewidth}{LRRRR}\toprule
指标 & 问题2 & 问题3 & 问题4-2 & 问题4-3\\\midrule
计划费用/万元 & 1356.38 & 1373.37 & 1439.75 & 1458.54\\
调整费用/万元 & 0.00 & 23.59 & 0.00 & 22.64\\
紧急费用/万元 & 49.48 & 7.31 & 47.71 & 6.01\\
总费用/万元 & 1405.86 & 1404.27 & 1487.45 & 1487.20\\
紧急电量/万kWh & 7.99 & 1.27 & 7.51 & 0.98\\
发生紧急购电天数 & 114.00 & 86.00 & 102.00 & 82.00\\
日费用CVaR95/万元 & 7.08 & 6.64 & 8.17 & 7.76\\
最高日费用/万元 & 9.24 & 6.96 & 10.00 & 8.14\\
2月1日期初电量/kWh & 8875.19 & 8814.60 & 9251.54 & 9030.40\\
12月31日期末电量/kWh & 6164.80 & 6247.51 & 6162.13 & 6312.01\\
\bottomrule\end{tabularx}\end{table}

\subsection{问题一的确定性调度结果}

问题一的计划购电量为59482.70 kWh，购电费用为35126.95元。无储能基准在每个时段按$[L_t-R_t]_+\Delta t$购电，费用为48052.05元，因此节约12925.10元，节费率为26.90\%。全天电网侧充电量20740.67 kWh、放电量16799.94 kWh，满足$0.9\sum_tc_t=\sum_td_t/0.9$；日首和日末储电量均为6000 kWh。该结果在给定预测视为确定输入的条件下成立，不能作为未知实际出力环境的完美预测基准直接与后续方案混比。

图\ref{fig:q1}呈现两段主要储能转移：低价夜间回补电量以支持早间负荷，白天吸收光伏与低价电量以支持晚间高价负荷。早间电池没有机械地放至最低值，晚峰附近则接近物理下限，体现后续光伏补能机会与价格机会成本的共同作用。晚间结束后恢复至首态，使节费不依赖透支初始库存。10分钟尺度电价的小幅变化会导致最优充电时段出现离散跳变，题目未给出爬坡或启停约束，故不额外平滑原始最优计划。

\begin{table}[htbp]\centering\small\caption{问题1指定时段与全天购电结果（电量单位：kWh）}\label{tab:q1-purchase}
\begin{tabularx}{\linewidth}{lRlRlR}
\toprule
时间段 & 购电量 & 时间段 & 购电量 & 时间段 & 购电量\\
\midrule
10:00--10:10 & 0.00 & 12:00--12:10 & 480.41 & 14:00--14:10 & 0.00\\
16:00--16:10 & 445.43 & 18:00--18:10 & 531.89 & 20:00--20:10 & 0.00\\
\midrule
全天购电量 & 59482.70 & \multicolumn{2}{c}{全天购电费（元）} & \multicolumn{2}{r}{35126.95}\\
\bottomrule
\end{tabularx}\end{table}
\begin{table}[htbp]\centering\small\caption{问题1储能充放电及首末电量（单位：kWh）}\label{tab:q1-battery}
\begin{tabularx}{\linewidth}{lRRlRR}
\toprule
时间段 & 充电量 & 放电量 & 时间段 & 充电量 & 放电量\\
\midrule
00:00--04:00 & 4500.00 & 0.00 & 04:00--08:00 & 833.33 & 6365.84\\
08:00--12:00 & 4787.96 & 1703.00 & 12:00--16:00 & 5286.04 & 91.10\\
16:00--20:00 & 0.00 & 5780.13 & 20:00--24:00 & 5333.33 & 2859.87\\
\midrule
0:00储电量 & \multicolumn{2}{r}{6000.00} & 24:00储电量 & \multicolumn{2}{r}{6000.00}\\
\bottomrule
\end{tabularx}\end{table}
\begin{figure}[htbp]
\centering
\includegraphics[width=0.94\linewidth]{figures/results/q1_dispatch.pdf}
\caption{问题一购电、充放电及储电量与电价的对应关系}
\label{fig:q1}
\end{figure}
\FloatBarrier

\subsection{问题二的风险与储能共同作用}

问题二评价期计划费用为13563772.88元，紧急费用为494809.59元，总费用14058582.47元。紧急购电量79905.90 kWh，涉及114天的550个原始10分钟时段。紧急补购后逐时满足负荷，因此这里的“紧急购电发生”表示额外昂贵补救，不能解释为失供电。

表\ref{tab:ablation}使用相同数据、评价范围与结算方式，比较无储能场景采购、点预测LP、场景期望成本、CVaR增强和参考备用增强。点预测LP费用为15918915.58元，紧急量达到614871.37 kWh，说明仅围绕均值进行精细套利，可能放大实际残差对应的五倍补购代价。主策略较该基准节省11.69\%，紧急量降低约87.00\%。无储能基准虽然采用场景采购，仍须支付18627630.00元，表明误差缓冲与电量跨时转移均具有价值。

场景期望成本模型已将费用降至14063286.32元，主策略为14058582.47元，两者总费用接近。与仅加入CVaR的14056083.78元相比，再加入参考备用多支付2498.70元，同时紧急量由90170.41降至79905.90 kWh，下降11.38\%。因此备用机制提供的是应急需求与费用之间的选择，不能把其作用夸大为无条件降低费用。

\begin{table}[htbp]\centering\small\caption{问题2的逐层对照（评价期334天）}
\label{tab:ablation}
\begin{tabularx}{\linewidth}{LRRRR}\toprule
策略 & 总费用/万元 & 紧急量/kWh & 日费CVaR95/元 & 年末电量/kWh\\\midrule
无储能+场景采购 & 1862.76 & 299833.62 & 80408.44 & 6000.00\\
点预测LP & 1591.89 & 614871.37 & 104699.88 & 5899.90\\
场景期望成本 & 1406.33 & 116065.71 & 72318.87 & 6150.47\\
场景+CVaR & 1405.61 & 90170.40 & 71037.47 & 6164.44\\
场景+CVaR+参考备用 & 1405.86 & 79905.90 & 70757.81 & 6164.80\\
\bottomrule\end{tabularx}\end{table}
\begin{table}[htbp]\centering\small
\caption{问题2指定日期购电结果（沿用题目表1布局，电量单位：kWh）}
\label{tab:q2-purchase}
\par\smallskip\textbf{2025-03-20}\par\smallskip
\begin{tabularx}{\linewidth}{lRlRlR}
\toprule
时间段 & 购电量 & 时间段 & 购电量 & 时间段 & 购电量\\
\midrule
10:00--10:10 & 0.00 & 12:00--12:10 & 634.67 & 14:00--14:10 & 0.00\\
16:00--16:10 & 485.00 & 18:00--18:10 & 692.70 & 20:00--20:10 & 0.00\\
\midrule
全天购电量 & 69564.03 & \multicolumn{2}{c}{全天购电费（元）} & \multicolumn{2}{r}{42461.59}\\
\bottomrule
\end{tabularx}
\par\smallskip\textbf{2025-06-21}\par\smallskip
\begin{tabularx}{\linewidth}{lRlRlR}
\toprule
时间段 & 购电量 & 时间段 & 购电量 & 时间段 & 购电量\\
\midrule
10:00--10:10 & 0.00 & 12:00--12:10 & 0.00 & 14:00--14:10 & 0.00\\
16:00--16:10 & 230.93 & 18:00--18:10 & 399.79 & 20:00--20:10 & 0.00\\
\midrule
全天购电量 & 40664.38 & \multicolumn{2}{c}{全天购电费（元）} & \multicolumn{2}{r}{24081.64}\\
\bottomrule
\end{tabularx}
\par\smallskip\textbf{2025-09-23}\par\smallskip
\begin{tabularx}{\linewidth}{lRlRlR}
\toprule
时间段 & 购电量 & 时间段 & 购电量 & 时间段 & 购电量\\
\midrule
10:00--10:10 & 0.00 & 12:00--12:10 & 443.93 & 14:00--14:10 & 0.00\\
16:00--16:10 & 535.12 & 18:00--18:10 & 762.12 & 20:00--20:10 & 17.80\\
\midrule
全天购电量 & 69074.66 & \multicolumn{2}{c}{全天购电费（元）} & \multicolumn{2}{r}{46508.63}\\
\bottomrule
\end{tabularx}
\par\smallskip\textbf{2025-12-21}\par\smallskip
\begin{tabularx}{\linewidth}{lRlRlR}
\toprule
时间段 & 购电量 & 时间段 & 购电量 & 时间段 & 购电量\\
\midrule
10:00--10:10 & 0.00 & 12:00--12:10 & 997.23 & 14:00--14:10 & 0.00\\
16:00--16:10 & 879.43 & 18:00--18:10 & 727.92 & 20:00--20:10 & 0.00\\
\midrule
全天购电量 & 99320.30 & \multicolumn{2}{c}{全天购电费（元）} & \multicolumn{2}{r}{64264.49}\\
\bottomrule
\end{tabularx}
\par\smallskip\parbox{\linewidth}{\footnotesize 注：问题2与4-2列示计划购电量；问题3与4-3列示最终生效的调整后常规购电量。全天购电量不含紧急购电，全天购电费含计划、调整与紧急费用。}
\end{table}

\begin{table}[htbp]\centering\small
\caption{问题2指定日期储能结果（沿用题目表2布局，单位：kWh）}
\label{tab:q2-battery}
\par\smallskip\textbf{2025-03-20}\par\smallskip
\begin{tabularx}{\linewidth}{lRRlRR}
\toprule
时间段 & 充电量 & 放电量 & 时间段 & 充电量 & 放电量\\
\midrule
00:00--04:00 & 3848.92 & 167.57 & 04:00--08:00 & 1282.09 & 4456.53\\
08:00--12:00 & 6288.36 & 2777.38 & 12:00--16:00 & 2581.36 & 417.19\\
16:00--20:00 & 260.42 & 5745.65 & 20:00--24:00 & 5531.76 & 2978.49\\
\midrule
0:00储电量 & \multicolumn{2}{r}{6707.60} & 24:00储电量 & \multicolumn{2}{r}{6140.32}\\
\bottomrule
\end{tabularx}
\par\smallskip\textbf{2025-06-21}\par\smallskip
\begin{tabularx}{\linewidth}{lRRlRR}
\toprule
时间段 & 充电量 & 放电量 & 时间段 & 充电量 & 放电量\\
\midrule
00:00--04:00 & 1059.25 & 14.93 & 04:00--08:00 & 535.52 & 2355.82\\
08:00--12:00 & 4639.89 & 0.00 & 12:00--16:00 & 0.00 & 0.00\\
16:00--20:00 & 3.17 & 4757.93 & 20:00--24:00 & 6031.16 & 2485.90\\
\midrule
0:00储电量 & \multicolumn{2}{r}{7822.98} & 24:00储电量 & \multicolumn{2}{r}{8182.20}\\
\bottomrule
\end{tabularx}
\par\smallskip\textbf{2025-09-23}\par\smallskip
\begin{tabularx}{\linewidth}{lRRlRR}
\toprule
时间段 & 充电量 & 放电量 & 时间段 & 充电量 & 放电量\\
\midrule
00:00--04:00 & 4672.35 & 162.52 & 04:00--08:00 & 798.33 & 5611.11\\
08:00--12:00 & 5607.93 & 1896.62 & 12:00--16:00 & 3859.91 & 594.11\\
16:00--20:00 & 200.18 & 5819.81 & 20:00--24:00 & 5559.88 & 2692.31\\
\midrule
0:00储电量 & \multicolumn{2}{r}{6173.75} & 24:00储电量 & \multicolumn{2}{r}{6161.95}\\
\bottomrule
\end{tabularx}
\par\smallskip\textbf{2025-12-21}\par\smallskip
\begin{tabularx}{\linewidth}{lRRlRR}
\toprule
时间段 & 充电量 & 放电量 & 时间段 & 充电量 & 放电量\\
\midrule
00:00--04:00 & 4432.44 & 92.00 & 04:00--08:00 & 1066.23 & 1424.56\\
08:00--12:00 & 6219.38 & 6304.75 & 12:00--16:00 & 5284.59 & 2220.11\\
16:00--20:00 & 7.12 & 5579.45 & 20:00--24:00 & 5642.68 & 2847.32\\
\midrule
0:00储电量 & \multicolumn{2}{r}{6654.76} & 24:00储电量 & \multicolumn{2}{r}{6521.75}\\
\bottomrule
\end{tabularx}
\end{table}

\begin{table}[htbp]\centering\small
\caption{问题2指定日期紧急购电（题目表3布局，单位：kWh）}
\label{tab:q2-emergency}
\setlength{\tabcolsep}{3pt}
\begin{tabularx}{\linewidth}{lRlRlRlR}\toprule
\multicolumn{2}{c}{2025-03-20} & \multicolumn{2}{c}{2025-06-21} & \multicolumn{2}{c}{2025-09-23} & \multicolumn{2}{c}{2025-12-21}\\
时间段 & 购电量 & 时间段 & 购电量 & 时间段 & 购电量 & 时间段 & 购电量\\\midrule
无 & 0.00 & 无 & 0.00 & 20:20--21:10 & 562.34 & 无 & 0.00\\
\midrule
合计 & 0.00 & 合计 & 0.00 & 合计 & 562.34 & 合计 & 0.00\\
\bottomrule\end{tabularx}\end{table}

四个指定日期中，9月23日20:20至21:10需紧急购电562.34 kWh，其余三个指定日期无紧急补购。这并不意味着全年其他日期没有误差风险。表\ref{tab:q2-battery}中的四小时充、放电量可以同时为正，因为它们分别累加该区间内不同10分钟时段的动作，并非同一时刻同时充放电。
\FloatBarrier

\subsection{问题三的滚动修正及预报价值}

问题三的计划费用为13733730.06元，调整费用为235908.13元，紧急费用为73088.34元，总费用14042726.52元。紧急电量下降至12685.74 kWh，涉及86天、276个10分钟时段。相对问题二，总费用仅降低0.11\%，紧急量降低84.12\%；由于两问零时光伏信息源不同，该横向差异不宜全部归为滚动优化的净贡献。

为单独考察更新机制，表\ref{tab:updates}固定问题三的零时光伏预报来源及其他参数，依次允许6时、12时和18时加入更新。相对全天不更新，四次更新策略节约168834.48元，即1.19\%，紧急电量降低83.03\%。各策略从相同1月1日初态出发，保持自身运行历史，因而这是全年可执行策略的对照，不是强制每次更新前电量相同的孤立预测实验。

6时更新首先将紧急量由74740.62 kWh降至28083.40 kWh；进一步引入12时更新，费用再降低94270.14元；加入18时更新，费用再降低52353.94元、紧急量由26491.20降至12685.74 kWh。18时以后光伏预报误差的绝对改善较小，但此时储电状态、负荷偏差与晚间高价购电机会均已发生变化，故重新优化仍有收益。这些收益不能只归因于某一份光伏预报。

\begin{table}[htbp]\centering\small\caption{问题3预报更新与反馈频率的对照}
\label{tab:updates}
\begin{tabularx}{\linewidth}{LRRRR}\toprule
策略 & 总费用/万元 & 紧急量/kWh & 日费CVaR95/元 & 年末电量/kWh\\\midrule
仅0时 & 1421.16 & 74740.62 & 70159.55 & 6247.51\\
0、6时 & 1418.94 & 28083.40 & 66393.05 & 6247.51\\
0、6、12时 & 1409.51 & 26491.20 & 66457.03 & 6247.51\\
0、6、12、18时 & 1404.27 & 12685.74 & 66381.40 & 6247.51\\
\bottomrule\end{tabularx}\end{table}
\begin{figure}[htbp]
\centering
\includegraphics[width=0.98\linewidth]{figures/results/update_value.pdf}
\caption{固定零时信息源后，增加规定更新时点的费用与紧急量变化}
\label{fig:update_value}
\end{figure}

据此，在题目给定的免费预报与调度条件下，建议使用6、12、18时的更新机会。若“其他时刻”指四个规定时点之外，则附件没有提供相应预报，无法给出增加到每小时或每10分钟更新后的实测费用结论。可在新增信息可获得后，以“期望减少的紧急费用减去新增调整费”与获取预报、运行优化器的成本比较，再决定是否增设时点，不能用全天真实光伏代替不存在的高频预报。

\begin{table}[htbp]\centering\small
\caption{问题3指定日期购电结果（沿用题目表1布局，电量单位：kWh）}
\label{tab:q3-purchase}
\par\smallskip\textbf{2025-03-20}\par\smallskip
\begin{tabularx}{\linewidth}{lRlRlR}
\toprule
时间段 & 购电量 & 时间段 & 购电量 & 时间段 & 购电量\\
\midrule
10:00--10:10 & 0.00 & 12:00--12:10 & 644.57 & 14:00--14:10 & 0.00\\
16:00--16:10 & 663.46 & 18:00--18:10 & 694.45 & 20:00--20:10 & 0.00\\
\midrule
全天购电量 & 72034.34 & \multicolumn{2}{c}{全天购电费（元）} & \multicolumn{2}{r}{44213.75}\\
\bottomrule
\end{tabularx}
\par\smallskip\textbf{2025-06-21}\par\smallskip
\begin{tabularx}{\linewidth}{lRlRlR}
\toprule
时间段 & 购电量 & 时间段 & 购电量 & 时间段 & 购电量\\
\midrule
10:00--10:10 & 0.00 & 12:00--12:10 & 0.00 & 14:00--14:10 & 0.00\\
16:00--16:10 & 216.64 & 18:00--18:10 & 321.63 & 20:00--20:10 & 0.00\\
\midrule
全天购电量 & 37297.18 & \multicolumn{2}{c}{全天购电费（元）} & \multicolumn{2}{r}{22283.44}\\
\bottomrule
\end{tabularx}
\par\smallskip\textbf{2025-09-23}\par\smallskip
\begin{tabularx}{\linewidth}{lRlRlR}
\toprule
时间段 & 购电量 & 时间段 & 购电量 & 时间段 & 购电量\\
\midrule
10:00--10:10 & 0.00 & 12:00--12:10 & 467.76 & 14:00--14:10 & 0.00\\
16:00--16:10 & 590.04 & 18:00--18:10 & 786.18 & 20:00--20:10 & 5.57\\
\midrule
全天购电量 & 73151.44 & \multicolumn{2}{c}{全天购电费（元）} & \multicolumn{2}{r}{45573.45}\\
\bottomrule
\end{tabularx}
\par\smallskip\textbf{2025-12-21}\par\smallskip
\begin{tabularx}{\linewidth}{lRlRlR}
\toprule
时间段 & 购电量 & 时间段 & 购电量 & 时间段 & 购电量\\
\midrule
10:00--10:10 & 0.00 & 12:00--12:10 & 992.64 & 14:00--14:10 & 0.00\\
16:00--16:10 & 897.85 & 18:00--18:10 & 724.70 & 20:00--20:10 & 0.00\\
\midrule
全天购电量 & 99253.15 & \multicolumn{2}{c}{全天购电费（元）} & \multicolumn{2}{r}{64091.54}\\
\bottomrule
\end{tabularx}
\par\smallskip\parbox{\linewidth}{\footnotesize 注：问题2与4-2列示计划购电量；问题3与4-3列示最终生效的调整后常规购电量。全天购电量不含紧急购电，全天购电费含计划、调整与紧急费用。}
\end{table}

\begin{table}[htbp]\centering\small
\caption{问题3指定日期储能结果（沿用题目表2布局，单位：kWh）}
\label{tab:q3-battery}
\par\smallskip\textbf{2025-03-20}\par\smallskip
\begin{tabularx}{\linewidth}{lRRlRR}
\toprule
时间段 & 充电量 & 放电量 & 时间段 & 充电量 & 放电量\\
\midrule
00:00--04:00 & 3970.31 & 143.39 & 04:00--08:00 & 900.49 & 5889.37\\
08:00--12:00 & 5942.16 & 2777.38 & 12:00--16:00 & 4959.17 & 254.72\\
16:00--20:00 & 17.80 & 5636.88 & 20:00--24:00 & 5546.09 & 2972.18\\
\midrule
0:00储电量 & \multicolumn{2}{r}{6677.15} & 24:00储电量 & \multicolumn{2}{r}{6241.87}\\
\bottomrule
\end{tabularx}
\par\smallskip\textbf{2025-06-21}\par\smallskip
\begin{tabularx}{\linewidth}{lRRlRR}
\toprule
时间段 & 充电量 & 放电量 & 时间段 & 充电量 & 放电量\\
\midrule
00:00--04:00 & 852.37 & 571.84 & 04:00--08:00 & 524.05 & 4390.39\\
08:00--12:00 & 8060.49 & 0.00 & 12:00--16:00 & 0.00 & 0.00\\
16:00--20:00 & 99.66 & 4836.09 & 20:00--24:00 & 6004.38 & 2485.90\\
\midrule
0:00储电量 & \multicolumn{2}{r}{7820.38} & 24:00储电量 & \multicolumn{2}{r}{8158.10}\\
\bottomrule
\end{tabularx}
\par\smallskip\textbf{2025-09-23}\par\smallskip
\begin{tabularx}{\linewidth}{lRRlRR}
\toprule
时间段 & 充电量 & 放电量 & 时间段 & 充电量 & 放电量\\
\midrule
00:00--04:00 & 4299.55 & 80.35 & 04:00--08:00 & 823.04 & 5847.79\\
08:00--12:00 & 6293.08 & 1896.62 & 12:00--16:00 & 3657.67 & 413.77\\
16:00--20:00 & 131.83 & 5662.28 & 20:00--24:00 & 5678.17 & 3091.35\\
\midrule
0:00储电量 & \multicolumn{2}{r}{6376.64} & 24:00储电量 & \multicolumn{2}{r}{6291.47}\\
\bottomrule
\end{tabularx}
\par\smallskip\textbf{2025-12-21}\par\smallskip
\begin{tabularx}{\linewidth}{lRRlRR}
\toprule
时间段 & 充电量 & 放电量 & 时间段 & 充电量 & 放电量\\
\midrule
00:00--04:00 & 4544.08 & 36.25 & 04:00--08:00 & 961.65 & 1476.18\\
08:00--12:00 & 6423.45 & 6283.79 & 12:00--16:00 & 5219.32 & 2316.39\\
16:00--20:00 & 6.05 & 5572.83 & 20:00--24:00 & 5660.24 & 2842.53\\
\midrule
0:00储电量 & \multicolumn{2}{r}{6602.57} & 24:00储电量 & \multicolumn{2}{r}{6549.26}\\
\bottomrule
\end{tabularx}
\end{table}

\begin{table}[htbp]\centering\small
\caption{问题3指定日期紧急购电（题目表3布局，单位：kWh）}
\label{tab:q3-emergency}
\setlength{\tabcolsep}{3pt}
\begin{tabularx}{\linewidth}{lRlRlRlR}\toprule
\multicolumn{2}{c}{2025-03-20} & \multicolumn{2}{c}{2025-06-21} & \multicolumn{2}{c}{2025-09-23} & \multicolumn{2}{c}{2025-12-21}\\
时间段 & 购电量 & 时间段 & 购电量 & 时间段 & 购电量 & 时间段 & 购电量\\\midrule
无 & 0.00 & 无 & 0.00 & 20:20--20:30 & 6.38 & 无 & 0.00\\
\midrule
合计 & 0.00 & 合计 & 0.00 & 合计 & 6.38 & 合计 & 0.00\\
\bottomrule\end{tabularx}\end{table}
\begin{table}[htbp]\centering\small
\caption{问题3指定日期计划、调整与费用分解}
\begin{tabularx}{\linewidth}{lRRRRR}\toprule
日期 & 计划量/kWh & 生效量/kWh & 计划费/元 & 调整费/元 & 紧急费/元\\\midrule
2025-03-20 & 71922.07 & 72034.34 & 44073.93 & 139.82 & 0.00\\
2025-06-21 & 37292.62 & 37297.18 & 22276.77 & 6.67 & 0.00\\
2025-09-23 & 73151.44 & 73151.44 & 45530.46 & 0.00 & 43.00\\
2025-12-21 & 99253.15 & 99253.15 & 64091.54 & 0.00 & 0.00\\
\bottomrule\end{tabularx}\end{table}

图\ref{fig:seasons}显示四个指定日期的零时计划与最终常规购电；多数时段两条曲线重合，只有新增购电必要时才发生修订。9月23日仍有20:20至20:30共6.38 kWh紧急购电，相比问题二该日562.34 kWh显著减少，但不能声称完全消除了应急补救。

\begin{figure}[htbp]
\centering
\includegraphics[width=0.96\linewidth]{figures/results/q3_seasons.pdf}
\caption{问题三四个指定日期的计划与执行；虚线表示6、12、18时更新节点}
\label{fig:seasons}
\end{figure}
\FloatBarrier

\subsection{问题四的波动电价结果}

在实际波动电价下，问题四对应问题二的策略总费用为14874522.48元，其中紧急费用477058.06元，紧急电量75130.77 kWh；对应问题三的滚动策略总费用为14871979.24元，其中调整费226405.77元、紧急费60145.05元，紧急电量9795.55 kWh。两者常规计划费用分别为14397464.42元和14585428.42元。

滚动策略节约2543.24元，仅占静态费用的0.017\%，但紧急电量降低86.96\%。这种结果符合题目收费结构：提早锁定较多电量和增加调整交易，能够减少五倍紧急购电，然而常规采购与调整成本抵消了大部分节省。日费用经验CVaR95由81744.59元降至77615.52元，最高日费用由100034.45元降至81371.51元，说明在本年度数据中更明显的收益体现在高费用日期与应急依赖，而非平均总费用的大幅降低。

\begin{table}[htbp]\centering\small
\caption{问题4-2指定日期购电结果（沿用题目表1布局，电量单位：kWh）}
\label{tab:q4_2-purchase}
\par\smallskip\textbf{2025-03-20}\par\smallskip
\begin{tabularx}{\linewidth}{lRlRlR}
\toprule
时间段 & 购电量 & 时间段 & 购电量 & 时间段 & 购电量\\
\midrule
10:00--10:10 & 0.00 & 12:00--12:10 & 634.67 & 14:00--14:10 & 0.00\\
16:00--16:10 & 501.72 & 18:00--18:10 & 0.00 & 20:00--20:10 & 735.49\\
\midrule
全天购电量 & 69863.61 & \multicolumn{2}{c}{全天购电费（元）} & \multicolumn{2}{r}{43714.39}\\
\bottomrule
\end{tabularx}
\par\smallskip\textbf{2025-06-21}\par\smallskip
\begin{tabularx}{\linewidth}{lRlRlR}
\toprule
时间段 & 购电量 & 时间段 & 购电量 & 时间段 & 购电量\\
\midrule
10:00--10:10 & 0.00 & 12:00--12:10 & 0.00 & 14:00--14:10 & 0.00\\
16:00--16:10 & 232.58 & 18:00--18:10 & 533.11 & 20:00--20:10 & 0.00\\
\midrule
全天购电量 & 42111.36 & \multicolumn{2}{c}{全天购电费（元）} & \multicolumn{2}{r}{21455.96}\\
\bottomrule
\end{tabularx}
\par\smallskip\textbf{2025-09-23}\par\smallskip
\begin{tabularx}{\linewidth}{lRlRlR}
\toprule
时间段 & 购电量 & 时间段 & 购电量 & 时间段 & 购电量\\
\midrule
10:00--10:10 & 0.00 & 12:00--12:10 & 588.49 & 14:00--14:10 & 0.00\\
16:00--16:10 & 572.26 & 18:00--18:10 & 0.00 & 20:00--20:10 & 17.80\\
\midrule
全天购电量 & 68805.32 & \multicolumn{2}{c}{全天购电费（元）} & \multicolumn{2}{r}{47789.54}\\
\bottomrule
\end{tabularx}
\par\smallskip\textbf{2025-12-21}\par\smallskip
\begin{tabularx}{\linewidth}{lRlRlR}
\toprule
时间段 & 购电量 & 时间段 & 购电量 & 时间段 & 购电量\\
\midrule
10:00--10:10 & 0.00 & 12:00--12:10 & 987.41 & 14:00--14:10 & 0.00\\
16:00--16:10 & 879.43 & 18:00--18:10 & 731.06 & 20:00--20:10 & 0.00\\
\midrule
全天购电量 & 99329.95 & \multicolumn{2}{c}{全天购电费（元）} & \multicolumn{2}{r}{75678.23}\\
\bottomrule
\end{tabularx}
\par\smallskip\parbox{\linewidth}{\footnotesize 注：问题2与4-2列示计划购电量；问题3与4-3列示最终生效的调整后常规购电量。全天购电量不含紧急购电，全天购电费含计划、调整与紧急费用。}
\end{table}

\begin{table}[htbp]\centering\small
\caption{问题4-2指定日期储能结果（沿用题目表2布局，单位：kWh）}
\label{tab:q4_2-battery}
\par\smallskip\textbf{2025-03-20}\par\smallskip
\begin{tabularx}{\linewidth}{lRRlRR}
\toprule
时间段 & 充电量 & 放电量 & 时间段 & 充电量 & 放电量\\
\midrule
00:00--04:00 & 2173.72 & 175.05 & 04:00--08:00 & 3492.41 & 5061.23\\
08:00--12:00 & 6291.46 & 2777.38 & 12:00--16:00 & 2694.62 & 356.34\\
16:00--20:00 & 191.41 & 6344.65 & 20:00--24:00 & 5580.52 & 2211.32\\
\midrule
0:00储电量 & \multicolumn{2}{r}{6800.10} & 24:00储电量 & \multicolumn{2}{r}{6375.20}\\
\bottomrule
\end{tabularx}
\par\smallskip\textbf{2025-06-21}\par\smallskip
\begin{tabularx}{\linewidth}{lRRlRR}
\toprule
时间段 & 充电量 & 放电量 & 时间段 & 充电量 & 放电量\\
\midrule
00:00--04:00 & 779.63 & 3116.33 & 04:00--08:00 & 2976.61 & 1719.95\\
08:00--12:00 & 5636.01 & 0.00 & 12:00--16:00 & 0.00 & 0.00\\
16:00--20:00 & 0.00 & 4809.67 & 20:00--24:00 & 5846.76 & 2485.90\\
\midrule
0:00储电量 & \multicolumn{2}{r}{7720.62} & 24:00储电量 & \multicolumn{2}{r}{7955.90}\\
\bottomrule
\end{tabularx}
\par\smallskip\textbf{2025-09-23}\par\smallskip
\begin{tabularx}{\linewidth}{lRRlRR}
\toprule
时间段 & 充电量 & 放电量 & 时间段 & 充电量 & 放电量\\
\midrule
00:00--04:00 & 4767.05 & 44.00 & 04:00--08:00 & 99.10 & 5848.24\\
08:00--12:00 & 5319.17 & 1896.62 & 12:00--16:00 & 4406.89 & 479.52\\
16:00--20:00 & 284.97 & 7326.64 & 20:00--24:00 & 5574.69 & 1301.33\\
\midrule
0:00储电量 & \multicolumn{2}{r}{6523.80} & 24:00储电量 & \multicolumn{2}{r}{6156.77}\\
\bottomrule
\end{tabularx}
\par\smallskip\textbf{2025-12-21}\par\smallskip
\begin{tabularx}{\linewidth}{lRRlRR}
\toprule
时间段 & 充电量 & 放电量 & 时间段 & 充电量 & 放电量\\
\midrule
00:00--04:00 & 4574.65 & 75.61 & 04:00--08:00 & 800.60 & 1392.25\\
08:00--12:00 & 5983.94 & 6334.74 & 12:00--16:00 & 5528.51 & 2220.11\\
16:00--20:00 & 7.12 & 5550.99 & 20:00--24:00 & 5691.21 & 2842.53\\
\midrule
0:00储电量 & \multicolumn{2}{r}{6737.42} & 24:00储电量 & \multicolumn{2}{r}{6602.37}\\
\bottomrule
\end{tabularx}
\end{table}

\begin{table}[htbp]\centering\small
\caption{问题4-2指定日期紧急购电（题目表3布局，单位：kWh）}
\label{tab:q4_2-emergency}
\setlength{\tabcolsep}{3pt}
\begin{tabularx}{\linewidth}{lRlRlRlR}\toprule
\multicolumn{2}{c}{2025-03-20} & \multicolumn{2}{c}{2025-06-21} & \multicolumn{2}{c}{2025-09-23} & \multicolumn{2}{c}{2025-12-21}\\
时间段 & 购电量 & 时间段 & 购电量 & 时间段 & 购电量 & 时间段 & 购电量\\\midrule
无 & 0.00 & 无 & 0.00 & 18:20--18:40 & 3.91 & 无 & 0.00\\
 &   &  &   & 20:30--21:00 & 481.54 &  &  \\
\midrule
合计 & 0.00 & 合计 & 0.00 & 合计 & 485.45 & 合计 & 0.00\\
\bottomrule\end{tabularx}\end{table}
\FloatBarrier
\begin{table}[htbp]\centering\small
\caption{问题4-3指定日期购电结果（沿用题目表1布局，电量单位：kWh）}
\label{tab:q4_3-purchase}
\par\smallskip\textbf{2025-03-20}\par\smallskip
\begin{tabularx}{\linewidth}{lRlRlR}
\toprule
时间段 & 购电量 & 时间段 & 购电量 & 时间段 & 购电量\\
\midrule
10:00--10:10 & 0.00 & 12:00--12:10 & 662.11 & 14:00--14:10 & 188.29\\
16:00--16:10 & 680.86 & 18:00--18:10 & 692.66 & 20:00--20:10 & 295.12\\
\midrule
全天购电量 & 71818.10 & \multicolumn{2}{c}{全天购电费（元）} & \multicolumn{2}{r}{45444.10}\\
\bottomrule
\end{tabularx}
\par\smallskip\textbf{2025-06-21}\par\smallskip
\begin{tabularx}{\linewidth}{lRlRlR}
\toprule
时间段 & 购电量 & 时间段 & 购电量 & 时间段 & 购电量\\
\midrule
10:00--10:10 & 0.00 & 12:00--12:10 & 0.00 & 14:00--14:10 & 0.00\\
16:00--16:10 & 216.84 & 18:00--18:10 & 514.10 & 20:00--20:10 & 0.00\\
\midrule
全天购电量 & 37850.75 & \multicolumn{2}{c}{全天购电费（元）} & \multicolumn{2}{r}{19741.04}\\
\bottomrule
\end{tabularx}
\par\smallskip\textbf{2025-09-23}\par\smallskip
\begin{tabularx}{\linewidth}{lRlRlR}
\toprule
时间段 & 购电量 & 时间段 & 购电量 & 时间段 & 购电量\\
\midrule
10:00--10:10 & 13.89 & 12:00--12:10 & 468.91 & 14:00--14:10 & 99.25\\
16:00--16:10 & 590.04 & 18:00--18:10 & 0.00 & 20:00--20:10 & 18.16\\
\midrule
全天购电量 & 73494.01 & \multicolumn{2}{c}{全天购电费（元）} & \multicolumn{2}{r}{48516.22}\\
\bottomrule
\end{tabularx}
\par\smallskip\textbf{2025-12-21}\par\smallskip
\begin{tabularx}{\linewidth}{lRlRlR}
\toprule
时间段 & 购电量 & 时间段 & 购电量 & 时间段 & 购电量\\
\midrule
10:00--10:10 & 0.00 & 12:00--12:10 & 1069.08 & 14:00--14:10 & 0.00\\
16:00--16:10 & 897.85 & 18:00--18:10 & 724.70 & 20:00--20:10 & 0.00\\
\midrule
全天购电量 & 99454.26 & \multicolumn{2}{c}{全天购电费（元）} & \multicolumn{2}{r}{75908.54}\\
\bottomrule
\end{tabularx}
\par\smallskip\parbox{\linewidth}{\footnotesize 注：问题2与4-2列示计划购电量；问题3与4-3列示最终生效的调整后常规购电量。全天购电量不含紧急购电，全天购电费含计划、调整与紧急费用。}
\end{table}

\begin{table}[htbp]\centering\small
\caption{问题4-3指定日期储能结果（沿用题目表2布局，单位：kWh）}
\label{tab:q4_3-battery}
\par\smallskip\textbf{2025-03-20}\par\smallskip
\begin{tabularx}{\linewidth}{lRRlRR}
\toprule
时间段 & 充电量 & 放电量 & 时间段 & 充电量 & 放电量\\
\midrule
00:00--04:00 & 4338.13 & 146.27 & 04:00--08:00 & 464.61 & 6007.86\\
08:00--12:00 & 6075.16 & 2777.38 & 12:00--16:00 & 4821.91 & 236.17\\
16:00--20:00 & 43.48 & 5767.09 & 20:00--24:00 & 5592.11 & 2652.48\\
\midrule
0:00储电量 & \multicolumn{2}{r}{6856.48} & 24:00储电量 & \multicolumn{2}{r}{6516.95}\\
\bottomrule
\end{tabularx}
\par\smallskip\textbf{2025-06-21}\par\smallskip
\begin{tabularx}{\linewidth}{lRRlRR}
\toprule
时间段 & 充电量 & 放电量 & 时间段 & 充电量 & 放电量\\
\midrule
00:00--04:00 & 606.72 & 3861.83 & 04:00--08:00 & 1440.68 & 1719.95\\
08:00--12:00 & 8112.31 & 0.00 & 12:00--16:00 & 0.00 & 0.00\\
16:00--20:00 & 0.00 & 4755.36 & 20:00--24:00 & 6799.43 & 3091.92\\
\midrule
0:00储电量 & \multicolumn{2}{r}{7858.24} & 24:00储电量 & \multicolumn{2}{r}{8200.29}\\
\bottomrule
\end{tabularx}
\par\smallskip\textbf{2025-09-23}\par\smallskip
\begin{tabularx}{\linewidth}{lRRlRR}
\toprule
时间段 & 充电量 & 放电量 & 时间段 & 充电量 & 放电量\\
\midrule
00:00--04:00 & 4735.94 & 72.91 & 04:00--08:00 & 367.84 & 5179.79\\
08:00--12:00 & 6326.29 & 1893.23 & 12:00--16:00 & 2296.37 & 201.04\\
16:00--20:00 & 194.97 & 6015.79 & 20:00--24:00 & 5659.23 & 2603.73\\
\midrule
0:00储电量 & \multicolumn{2}{r}{6609.50} & 24:00储电量 & \multicolumn{2}{r}{6491.54}\\
\bottomrule
\end{tabularx}
\par\smallskip\textbf{2025-12-21}\par\smallskip
\begin{tabularx}{\linewidth}{lRRlRR}
\toprule
时间段 & 充电量 & 放电量 & 时间段 & 充电量 & 放电量\\
\midrule
00:00--04:00 & 4651.42 & 90.02 & 04:00--08:00 & 0.00 & 825.69\\
08:00--12:00 & 6331.05 & 6029.43 & 12:00--16:00 & 4930.50 & 2266.73\\
16:00--20:00 & 0.00 & 5557.59 & 20:00--24:00 & 5703.18 & 2842.53\\
\midrule
0:00储电量 & \multicolumn{2}{r}{6713.74} & 24:00储电量 & \multicolumn{2}{r}{6599.39}\\
\bottomrule
\end{tabularx}
\end{table}

\begin{table}[htbp]\centering\small
\caption{问题4-3指定日期紧急购电（题目表3布局，单位：kWh）}
\label{tab:q4_3-emergency}
\setlength{\tabcolsep}{3pt}
\begin{tabularx}{\linewidth}{lRlRlRlR}\toprule
\multicolumn{2}{c}{2025-03-20} & \multicolumn{2}{c}{2025-06-21} & \multicolumn{2}{c}{2025-09-23} & \multicolumn{2}{c}{2025-12-21}\\
时间段 & 购电量 & 时间段 & 购电量 & 时间段 & 购电量 & 时间段 & 购电量\\\midrule
无 & 0.00 & 无 & 0.00 & 20:10--20:20 & 10.23 & 无 & 0.00\\
\midrule
合计 & 0.00 & 合计 & 0.00 & 合计 & 10.23 & 合计 & 0.00\\
\bottomrule\end{tabularx}\end{table}
\begin{table}[htbp]\centering\small
\caption{问题4-3指定日期计划、调整与费用分解}
\begin{tabularx}{\linewidth}{lRRRRR}\toprule
日期 & 计划量/kWh & 生效量/kWh & 计划费/元 & 调整费/元 & 紧急费/元\\\midrule
2025-03-20 & 71655.38 & 71818.10 & 45192.28 & 251.82 & 0.00\\
2025-06-21 & 37850.75 & 37850.75 & 19741.04 & 0.00 & 0.00\\
2025-09-23 & 73480.32 & 73494.01 & 48421.73 & 28.22 & 66.27\\
2025-12-21 & 99257.81 & 99454.26 & 75595.15 & 313.39 & 0.00\\
\bottomrule\end{tabularx}\end{table}

图\ref{fig:monthly}给出月度费用与紧急量。夏季及冬季的负荷、光伏和价格水平不同，固定电价与波动电价的结果会产生季节性差别。图中全部月份均为样本外顺序结算结果，历史窗口和风险参数不根据这些月份的评价值进行反向选择。

\begin{figure}[htbp]
\centering
\includegraphics[width=0.94\linewidth]{figures/results/monthly_results.pdf}
\caption{四组策略在评价期内的月度费用与紧急购电量}
\label{fig:monthly}
\end{figure}
\FloatBarrier

\section{参数敏感性与模型评价}

\subsection{备用、风险权重与终端库存}

主策略采用备用系数$\kappa=1$、风险基数$\lambda_0=0.1$。二者作为事先固定的策略偏好，未利用2月至12月的实际费用寻优。表\ref{tab:sensitivity}完整保留增强参数带来费用上升的情况，以判断机制的实际边界。

备用系数从0.5提高至1.5时，问题三紧急量由14208.94 kWh降至11688.95 kWh，总费用由14023535.12元升至14067387.71元。更多参考备用增加了缓冲空间，但也限制了廉价电量的灵活使用，并可能导致更多承诺购电剩余。若运行目标仅是最小化本样本实际电费，较大的备用系数不一定合适。

风险基数取0时，总费用13927283.23元、紧急量14609.38 kWh；主策略费用14042726.52元、紧急量12685.74 kWh。主策略以约11.54万元额外支出减少约1923.64 kWh紧急需求，而且其日总费用CVaR95也高于零风险权重情形。原因在于优化惩罚的是预测场景中的紧急费用尾部，实际总费用还包含常规采购和调整费，且经验场景存在有限样本误差。不能将对一个随机损失分量的风险约束，直接解释为对全部实现成本尾部的单调改进。若不对紧急需求额外赋予偏好，应优先考虑零风险权重的经济型方案；本文保留既定主参数以展示明确的风险偏好，而不在评价期事后挑选最佳参数。

\begin{table}[htbp]\centering\small\caption{问题3风险与终端参数敏感性}
\label{tab:sensitivity}
\begin{tabularx}{\linewidth}{LRRRR}\toprule
策略 & 总费用/万元 & 紧急量/kWh & 日费CVaR95/元 & 年末电量/kWh\\\midrule
主策略：$\kappa=1,\lambda_0=0.1$ & 1404.27 & 12685.74 & 66381.40 & 6247.51\\
备用系数$\kappa=0.5$ & 1402.35 & 14208.94 & 66336.84 & 6247.51\\
备用系数$\kappa=1.5$ & 1406.74 & 11688.95 & 66461.26 & 6247.51\\
风险基数$\lambda_0=0$ & 1392.73 & 14609.38 & 65683.92 & 6245.12\\
风险基数$\lambda_0=0.2$ & 1416.64 & 11312.96 & 66878.46 & 6285.62\\
取消终端偏离惩罚 & 1403.32 & 18350.10 & 66351.31 & 1637.69\\
下调退款解释 & 1369.08 & 13430.61 & 64275.62 & 6234.10\\
\bottomrule\end{tabularx}\end{table}

取消终端惩罚后，费用降至14033228.53元，但紧急量升至18350.10 kWh，年末储电量降至1637.69 kWh，明显低于主策略6247.51 kWh。由于结存能量不同，不能把9497.99元费用下降全部认定为运营效率提升。对终端机制的比较应同时报告期初、期末电量与实际紧急需求；软目标只能抑制有限时域内过度放空的倾向，不能保证全天实际末态恰等于参考值。

\subsection{下调购电是否退款的解释敏感性}

主结果严格采用式\eqref{eq:settlement}的“计划全额付费加调整相关费用”口径。另一种可能解释是，下调时取消原计划对应电费，再支付50\%违约费，于是其净影响为$-0.5p_t[g_t-a_t]_+$。为避免把语义差异伪装成算法优势，另以这一退款解释重算问题三，结果仅用于敏感性比较，不替换附件5主结果。

退款解释下，总费用为13690758.70元，最终常规购电量少于原始计划量，累计下调1223373.21 kWh，调整净费用为$-160832.93$元。与主口径的费用差异约35.20万元，明显大于主策略在若干消融对照中的改善。因此结算语义对结果具有实质影响。该敏感性说明所报数值依赖明示的结算假设，不能在不同退款规则之间直接比较策略优劣。

\subsection{模型特点与适用边界}

本文的主要特点是维持一致的物理、信息与结算结构。各问使用同一储电状态方程和执行器，区别仅在预测来源、场景维度与计划更新时间。整日联合残差保留相关误差，避免逐时独立抽样低估持续预测偏差；参考备用把预测误差转换为储能量级，便于解释其为何抑制过早放电。所有费用以实际执行和实际电价结算，使优化参考成本与真实支付结果能够分开比较。

物理可行性通过逐时能量守恒、储电边界、充放电上限、互斥与跨日连续性计算检查。执行器每次先使用物理可行储能，再紧急补足负荷，所以即使场景没有覆盖实际误差，也不会把未满足负荷隐藏为零购电。外网容量未受限制是这一性质的前提，若实际存在接口限额，需进一步引入网络约束和失供电代价。

仍有四方面限制。第一，模型仅使用一个年度，且最初只有一个月用于选择预测器，难以据此声称跨年份泛化。第二，两阶段模型固定各场景共享的储能参考轨迹，实际再由反馈规则修正，没有求解任意历史依赖控制策略。第三，滚动时域只延伸到当日末端，跨日峰谷变化仅通过软终端目标近似。第四，容量、效率和响应速度按题目常数处理，未引入寿命退化、爬坡、交易次数或真实电网容量，因此结果适用于题目给定边界，不能无修改用于实际设备运行。

\section{结论}

确定性条件下，线性储能联合调度以35126.95元完成单日供电，节费26.90\%；当预测误差存在时，按历史残差配置购电和储能参考轨迹，比单纯围绕均值套利显著减少五倍补购造成的费用。规定时点的滚动更新在固定电价下具有可观的应急需求削减效果；电价同时波动时，总费用改善较小，主要表现为较低的紧急电量和较低的高费用日期成本。

风险机制应作为可选择的偏好参数，不能预设其总能降低全部成本。以实际电费为唯一目标与希望更少依赖紧急购电，是两种应明确区分的运行偏好。完整的预测、求解、执行与结算代码以及全部334天策略随支撑材料保存，便于根据最终采用的偏好和结算解释复算。


\FloatBarrier
\CumcmAIUsed{建模方案整理、预测与调度程序生成、数据计算、文献检索核验、图表生成及论文初稿撰写}
% 文献元数据核验日期：2026-09-10。
% DOI 注册元数据与出版社页面的核验记录见 ../文献与模型对应.md。
% 本文件直接使用 % 文献元数据核验日期：2026-09-10。
% DOI 注册元数据与出版社页面的核验记录见 ../文献与模型对应.md。
% 本文件直接使用 % 文献元数据核验日期：2026-09-10。
% DOI 注册元数据与出版社页面的核验记录见 ../文献与模型对应.md。
% 本文件直接使用 \input{sections/references}，无需 BibTeX。
\begin{thebibliography}{99}

\bibitem{kong2016}
Kong X, Bai L, Hu Q, et al.
Day-ahead optimal scheduling method for grid-connected microgrid based on energy storage control strategy[J].
Journal of Modern Power Systems and Clean Energy, 2016, 4(4): 648--658.
DOI: \url{https://doi.org/10.1007/s40565-016-0245-0}.

\bibitem{silva2020}
Silva V A, Aoki A R, Lambert-Torres G.
Optimal day-ahead scheduling of microgrids with battery energy storage system[J].
Energies, 2020, 13(19): 5188.
DOI: \url{https://doi.org/10.3390/en13195188}.

\bibitem{abunima2022}
Abunima H, Park W-H, Glick M B, et al.
Two-stage stochastic optimization for operating a renewable-based microgrid[J].
Applied Energy, 2022, 325: 119848.
DOI: \url{https://doi.org/10.1016/j.apenergy.2022.119848}.

\bibitem{marino2018}
Marino C, Quddus M A, Marufuzzaman M, et al.
A chance-constrained two-stage stochastic programming model for reliable microgrid operations under power demand uncertainty[J].
Sustainable Energy, Grids and Networks, 2018, 13: 66--77.
DOI: \url{https://doi.org/10.1016/j.segan.2017.12.007}.

\bibitem{grover2018}
Grover-Silva E, Heleno M, Mashayekh S, et al.
A stochastic optimal power flow for scheduling flexible resources in microgrids operation[J].
Applied Energy, 2018, 229: 201--208.
DOI: \url{https://doi.org/10.1016/j.apenergy.2018.07.114}.

\bibitem{palma2013}
Palma-Behnke R, Benavides C, Lanas F, et al.
A microgrid energy management system based on the rolling horizon strategy[J].
IEEE Transactions on Smart Grid, 2013, 4(2): 996--1006.
DOI: \url{https://doi.org/10.1109/TSG.2012.2231440}.

\bibitem{silvente2015}
Silvente J, Kopanos G M, Pistikopoulos E N, et al.
A rolling horizon optimization framework for the simultaneous energy supply and demand planning in microgrids[J].
Applied Energy, 2015, 155: 485--501.
DOI: \url{https://doi.org/10.1016/j.apenergy.2015.05.090}.

\bibitem{li2016}
Li Z, Zang C, Zeng P, et al.
Combined two-stage stochastic programming and receding horizon control strategy for microgrid energy management considering uncertainty[J].
Energies, 2016, 9(7): 499.
DOI: \url{https://doi.org/10.3390/en9070499}.

\bibitem{jani2022}
Jani A, Karimi H, Jadid S.
Two-layer stochastic day-ahead and real-time energy management of networked microgrids considering integration of renewable energy resources[J].
Applied Energy, 2022, 323: 119630.
DOI: \url{https://doi.org/10.1016/j.apenergy.2022.119630}.

\bibitem{li2020}
Li J, Tan C, Ren Z, et al.
A two-stage optimal scheduling model of microgrid based on chance-constrained programming in spot markets[J].
Processes, 2020, 8(1): 107.
DOI: \url{https://doi.org/10.3390/pr8010107}.

\bibitem{rockafellar2000}
Rockafellar R T, Uryasev S.
Optimization of conditional value-at-risk[J].
The Journal of Risk, 2000, 2(3): 21--41.
DOI: \url{https://doi.org/10.21314/JOR.2000.038}.

\end{thebibliography}
，无需 BibTeX。
\begin{thebibliography}{99}

\bibitem{kong2016}
Kong X, Bai L, Hu Q, et al.
Day-ahead optimal scheduling method for grid-connected microgrid based on energy storage control strategy[J].
Journal of Modern Power Systems and Clean Energy, 2016, 4(4): 648--658.
DOI: \url{https://doi.org/10.1007/s40565-016-0245-0}.

\bibitem{silva2020}
Silva V A, Aoki A R, Lambert-Torres G.
Optimal day-ahead scheduling of microgrids with battery energy storage system[J].
Energies, 2020, 13(19): 5188.
DOI: \url{https://doi.org/10.3390/en13195188}.

\bibitem{abunima2022}
Abunima H, Park W-H, Glick M B, et al.
Two-stage stochastic optimization for operating a renewable-based microgrid[J].
Applied Energy, 2022, 325: 119848.
DOI: \url{https://doi.org/10.1016/j.apenergy.2022.119848}.

\bibitem{marino2018}
Marino C, Quddus M A, Marufuzzaman M, et al.
A chance-constrained two-stage stochastic programming model for reliable microgrid operations under power demand uncertainty[J].
Sustainable Energy, Grids and Networks, 2018, 13: 66--77.
DOI: \url{https://doi.org/10.1016/j.segan.2017.12.007}.

\bibitem{grover2018}
Grover-Silva E, Heleno M, Mashayekh S, et al.
A stochastic optimal power flow for scheduling flexible resources in microgrids operation[J].
Applied Energy, 2018, 229: 201--208.
DOI: \url{https://doi.org/10.1016/j.apenergy.2018.07.114}.

\bibitem{palma2013}
Palma-Behnke R, Benavides C, Lanas F, et al.
A microgrid energy management system based on the rolling horizon strategy[J].
IEEE Transactions on Smart Grid, 2013, 4(2): 996--1006.
DOI: \url{https://doi.org/10.1109/TSG.2012.2231440}.

\bibitem{silvente2015}
Silvente J, Kopanos G M, Pistikopoulos E N, et al.
A rolling horizon optimization framework for the simultaneous energy supply and demand planning in microgrids[J].
Applied Energy, 2015, 155: 485--501.
DOI: \url{https://doi.org/10.1016/j.apenergy.2015.05.090}.

\bibitem{li2016}
Li Z, Zang C, Zeng P, et al.
Combined two-stage stochastic programming and receding horizon control strategy for microgrid energy management considering uncertainty[J].
Energies, 2016, 9(7): 499.
DOI: \url{https://doi.org/10.3390/en9070499}.

\bibitem{jani2022}
Jani A, Karimi H, Jadid S.
Two-layer stochastic day-ahead and real-time energy management of networked microgrids considering integration of renewable energy resources[J].
Applied Energy, 2022, 323: 119630.
DOI: \url{https://doi.org/10.1016/j.apenergy.2022.119630}.

\bibitem{li2020}
Li J, Tan C, Ren Z, et al.
A two-stage optimal scheduling model of microgrid based on chance-constrained programming in spot markets[J].
Processes, 2020, 8(1): 107.
DOI: \url{https://doi.org/10.3390/pr8010107}.

\bibitem{rockafellar2000}
Rockafellar R T, Uryasev S.
Optimization of conditional value-at-risk[J].
The Journal of Risk, 2000, 2(3): 21--41.
DOI: \url{https://doi.org/10.21314/JOR.2000.038}.

\end{thebibliography}
，无需 BibTeX。
\begin{thebibliography}{99}

\bibitem{kong2016}
Kong X, Bai L, Hu Q, et al.
Day-ahead optimal scheduling method for grid-connected microgrid based on energy storage control strategy[J].
Journal of Modern Power Systems and Clean Energy, 2016, 4(4): 648--658.
DOI: \url{https://doi.org/10.1007/s40565-016-0245-0}.

\bibitem{silva2020}
Silva V A, Aoki A R, Lambert-Torres G.
Optimal day-ahead scheduling of microgrids with battery energy storage system[J].
Energies, 2020, 13(19): 5188.
DOI: \url{https://doi.org/10.3390/en13195188}.

\bibitem{abunima2022}
Abunima H, Park W-H, Glick M B, et al.
Two-stage stochastic optimization for operating a renewable-based microgrid[J].
Applied Energy, 2022, 325: 119848.
DOI: \url{https://doi.org/10.1016/j.apenergy.2022.119848}.

\bibitem{marino2018}
Marino C, Quddus M A, Marufuzzaman M, et al.
A chance-constrained two-stage stochastic programming model for reliable microgrid operations under power demand uncertainty[J].
Sustainable Energy, Grids and Networks, 2018, 13: 66--77.
DOI: \url{https://doi.org/10.1016/j.segan.2017.12.007}.

\bibitem{grover2018}
Grover-Silva E, Heleno M, Mashayekh S, et al.
A stochastic optimal power flow for scheduling flexible resources in microgrids operation[J].
Applied Energy, 2018, 229: 201--208.
DOI: \url{https://doi.org/10.1016/j.apenergy.2018.07.114}.

\bibitem{palma2013}
Palma-Behnke R, Benavides C, Lanas F, et al.
A microgrid energy management system based on the rolling horizon strategy[J].
IEEE Transactions on Smart Grid, 2013, 4(2): 996--1006.
DOI: \url{https://doi.org/10.1109/TSG.2012.2231440}.

\bibitem{silvente2015}
Silvente J, Kopanos G M, Pistikopoulos E N, et al.
A rolling horizon optimization framework for the simultaneous energy supply and demand planning in microgrids[J].
Applied Energy, 2015, 155: 485--501.
DOI: \url{https://doi.org/10.1016/j.apenergy.2015.05.090}.

\bibitem{li2016}
Li Z, Zang C, Zeng P, et al.
Combined two-stage stochastic programming and receding horizon control strategy for microgrid energy management considering uncertainty[J].
Energies, 2016, 9(7): 499.
DOI: \url{https://doi.org/10.3390/en9070499}.

\bibitem{jani2022}
Jani A, Karimi H, Jadid S.
Two-layer stochastic day-ahead and real-time energy management of networked microgrids considering integration of renewable energy resources[J].
Applied Energy, 2022, 323: 119630.
DOI: \url{https://doi.org/10.1016/j.apenergy.2022.119630}.

\bibitem{li2020}
Li J, Tan C, Ren Z, et al.
A two-stage optimal scheduling model of microgrid based on chance-constrained programming in spot markets[J].
Processes, 2020, 8(1): 107.
DOI: \url{https://doi.org/10.3390/pr8010107}.

\bibitem{rockafellar2000}
Rockafellar R T, Uryasev S.
Optimization of conditional value-at-risk[J].
The Journal of Risk, 2000, 2(3): 21--41.
DOI: \url{https://doi.org/10.21314/JOR.2000.038}.

\end{thebibliography}


\clearpage
\begin{appendices}
\section{支撑材料与可复现计算}

论文主文件为\texttt{main.tex}，原模板示例\texttt{example.tex}和类、样式文件保留。计算过程不依赖LaTeX；表格中的结果均由程序从已求解的数组生成，不手工改写求解值。支撑材料主要组成见表\ref{tab:files}。

\begin{table}[htbp]\centering\small
\caption{主要代码及结果文件}\label{tab:files}
\begin{tabularx}{\linewidth}{lL}\toprule
文件 & 内容\\\midrule
\texttt{code/data\_forecast.py} & 原始附件读取、因果预测、残差与数据统计\\
\texttt{code/optimize\_dispatch.py} & 统一LP、场景缩减、执行器与四组全年主策略\\
\texttt{code/run\_experiments.py} & 消融、更新时点、风险参数与退款解释实验\\
\texttt{code/make\_paper\_assets.py} & 指定日期表格、年度汇总、CSV和论文插图\\
\texttt{code/export\_workbooks.mjs} & 按原始模板填写附件5五个工作簿\\
\texttt{code/verify\_workbooks.py} & 输出与数组、表头、费用及日期的逐项一致性检查\\
\texttt{data/prepared.npz} & 数据、预测、各发布时间残差\\
\texttt{data/*\_full.npz} & 包含1月状态推进的全年逐时执行结果\\
\texttt{data/daily\_summary.csv} & 2月至12月常规量、紧急量、费用和库存汇总\\
\texttt{data/experiments.json} & 全部对照设置、结果与库存边界\\
\texttt{originals/} & 原附件5空模板及原论文模板代码备份\\
\bottomrule\end{tabularx}
\end{table}

各\texttt{*\_full.npz}中，\texttt{plan}为每日零时承诺量，\texttt{adjusted}为每个时段最终生效的常规购电量，\texttt{charge}与\texttt{discharge}为实际电网侧电量，\texttt{energy}包含一天145个边界状态，\texttt{emergency}为紧急量，\texttt{spill}为未使用的富余电量。\texttt{plan\_cost}、\texttt{adjustment\_cost}、\texttt{emergency\_cost}之和为\texttt{total\_cost}。未使用原始真实未来值重算并替换已执行的决策。

\section{原模板时间标签与费用填写说明}

附件5购电页首标签为“0:10--0:20”，末标签为“0:00--0:10+1”；按要求完整保留这些标签及顺序，数组与附件1、2、4的144个源数据列逐列对应。物理模型把源数据0:10解释为0:00--0:10的右端点，因此模板标签仅保留原始展示口径，不再据其字面区间将数据第二次移位。三种对应关系为：
\begin{center}\small
\begin{tabular}{lll}\toprule
源采样列 & 模型物理时段 & 保留的模板显示标签\\\midrule
第1列：0:10 & 0:00--0:10 & 0:10--0:20\\
第61列：10:10 & 10:00--10:10 & 10:10--10:20\\
第144列：0:00+1 & 23:50--24:00 & 0:00--0:10+1\\
\bottomrule\end{tabular}
\end{center}
上述约定保留源列次序与自然日能量守恒。正文指定时段表、四小时储能分段及紧急时段使用物理时钟，因此不能直接以购电页标签字符串在正文表中查找同一结果，应按上表映射源列。

计划购电页全天购电量为144段计划量之和，费用为计划付费部分；调整购电页填最终生效总量，不填相对计划的差额，全天费用为含计划、调整和紧急购电的实际总费。问题二与问题四对应问题二的工作簿没有调整页，其含紧急费的总费用见论文和\texttt{daily\_summary.csv}。充放电页按日填写六个四小时区间，并列日首和日末储电量。紧急页将相邻正紧急量的10分钟区间合并，跨日不合并，无紧急量的日期写“无”和0。

\section{执行与结算核心代码}

以下代码节选展示因果执行环节。完整程序及全部参数保存在独立源文件中，结果生成不依赖人工修改表格。
\begin{lstlisting}[language=Python,basicstyle=\ttfamily\footnotesize]
def step(grid, net, energy):
    cap, eta = 5000 / 6, 0.9
    if grid >= net:
        charge = min(grid - net, cap, (10800 - energy) / eta)
        discharge, emergency = 0.0, 0.0
        spill = grid - net - charge
    else:
        discharge = min(net - grid, cap, (energy - 1200) * eta)
        charge, spill = 0.0, 0.0
        emergency = net - grid - discharge
    next_energy = energy + eta * charge - discharge / eta
    return charge, discharge, next_energy, emergency, spill

plan_cost = sum(price * plan)
up = maximum(adjusted - plan, 0)
down = maximum(plan - adjusted, 0)
adjustment_cost = sum(price * (1.5 * up + 0.5 * down))
emergency_cost = sum(5 * price * emergency)
total_cost = plan_cost + adjustment_cost + emergency_cost
\end{lstlisting}

AI工具实际参与范围、主要提示和待人工处理事项另存\texttt{AI工具使用详情.tex}及说明文档；该材料不宣称已完成用户尚未实施的人工核验。

\end{appendices}
\end{document}
